\documentclass{article}
\PassOptionsToPackage{numbers,sort&compress}{natbib}
\usepackage[preprint]{neurips_2026}
\usepackage[T1]{fontenc}
\usepackage[utf8]{inputenc}
\usepackage{amsmath,amssymb,booktabs}
\usepackage{hyperref}
\hypersetup{hidelinks}
\title{Neural Memory and Interval Calibration}
\author{Chengxuan YUAN \quad Student ID: 50025581\\DSAA 2012, HKUST(GZ)}
\begin{document}
\maketitle

\section{Motivation and research questions}
Financial prediction intervals should expand when volatility rises and contract
after a temporary shock. Overall coverage can conceal slow adaptation during
these transitions. A neural forecaster estimates the current scale from past
returns, while an online calibrator uses historical normalized errors to adjust
interval width. These two uses of history may reinforce or offset each other.

The distinction is concrete: an observed return of magnitude two leaves a score
of two if its originally issued scale was one, but a score of one if that scale
was two. Even when today's scales agree, different past denominators can
produce different calibrated intervals. Input context length alone therefore
does not identify the model's effective learned dependence.

Neural conformal control, joint score--threshold adaptation, retrospective
residual updates, and strongly adaptive calibration already address changing
environments \citep{ncc,adapt,retro,saocp}. Training underspecification and
initialization stability in neural volatility forecasting also have prior
literature \citep{damour,harnet}. I will study the specific pathways through
which a fixed neural scale model interacts with an issued-score archive:

\textbf{RQ1:} How does the context-by-score-window interaction differ between
shock entry and recovery, across fitted neural models and shock durations?

\textbf{RQ2:} How do current-scale and historical score-denominator effects
combine to produce the interaction, and how does their decomposition depend
on the order of changes?

\textbf{RQ3:} Which patterns persist across independent training datasets and
initialization/sampling seeds, and what do validation response probes reveal
about the limits of a simple memory-based explanation?

\section{Approach and controlled design}
I train a GRU and a causal dilated TCN from raw past-return features using
Gaussian negative log-likelihood, Adam, and validation early stopping.
A fixed zero center isolates scale adaptation. For current context $a$,
archive context $b$, and score window $W$, the interval half-width is
\[
 h_t^{a,b,W}=s_t^a\,
 Q_{1-\alpha}\!\left(\left\{|y_i|/s_i^b\right\}_{i=t-W}^{t-1}\right),
 \qquad \alpha=0.1.
\]
Each scale uses only earlier observations; $Q$ selects past-score rank
$k=\lceil(W+1)(1-\alpha)\rceil$. The full crossing uses
$a,b\in\{16,64\}$ and $W\in\{32,128\}$: eight configurations per frozen fit.
The ordinary configuration has $a=b$; crossed contexts are diagnostic
interventions. Contexts share the same learned weights within each fit.
An issued historical denominator remains tied to its own past prediction.

Paired shock/no-shock trajectories share baseline innovations. I compare how
the effect of changing context varies with score-window length, then subtract
the same contrast on the matched no-shock path. Changing the current pathway
first or the archive pathway first gives two exact decompositions of this
interaction; I retain both. Their components are computational differences,
not unique causal mediation percentages, and their sign is not an absolute
model ranking.

\clearpage
\section{Evaluation, data, and validity}
Controlled simulations vary pulse duration and magnitude, Gaussian versus
heavy-tailed innovations, and latent stochastic volatility. True scales and
event positions are used only to construct paths and evaluate phases; they are
not forecasting inputs or tuning signals. Non-neural controls include EWMA and
training-fitted GARCH. Calibration controls include rolling quantiles, adaptive
feedback, and strongly adaptive online calibration \citep{saocp}.

The main metric is interval score,
$2h+(2/\alpha)(|y|-h)_+$; I also report coverage, width, and simulation-oracle
diagnostics. Entry and recovery are evaluated separately so that wider
intervals cannot count as improvement solely because coverage increases.
Score-window changes also alter finite quantile ranks, so they are not a pure
change in history length.

The existing replication uses three independent training/validation datasets,
five paired initialization/sampling seeds per dataset, and 128 paired test
paths per noise family. Seed indices are shared across datasets, so this is a
crossed design. Paired-path intervals conditional on saved models are reported
separately from retraining sensitivity; neither model--path combinations nor
shared duration prefixes are counted as independent replications.

Only generated data and public ECB USD/JPY/GBP/CHF reference rates per EUR are
used. The financial pilot uses daily log returns, trains through 2018,
validates in 2019--2020, and treats the already-inspected period from 2021 as
exploratory. Date-block comparisons retain cross-currency dependence.
New methods require fresh, separately frozen evaluation. Restricted
competition data are excluded.

\section{Expected contribution, feasibility, and plan}
The intended contribution is a reproducible mechanism study separating current
neural scale response from archived score denominators, with phase-specific
effects and explicit failure boundaries. Its value lies in explaining these
pathways rather than simply applying two networks or selecting the best score.
No new calibration algorithm or coverage theorem is assumed.

Exploration and the three-dataset replication are already complete on local
CPU. In the Gaussian reference case, the mean GRU-minus-TCN entry interaction
contrast retains its sign in all three new training datasets, while recovery
does not consistently replicate. The crossed entry sensitivity interval
includes zero, limiting architecture claims. Validation lag gradients,
older-history permutation, and pulse probes provide diagnostics; their
predictive usefulness is not yet established.

Following teaching-team feedback, I will consolidate the research question,
literature positioning, and mechanism interpretation in October. November
will focus on justified robustness checks and public financial analysis.
December 1--10 will cover the four-page report, slides, and 8--10-minute video,
ahead of the December 11 deadline. Deliverables include code, paired results,
diagnostic figures, and a clear account of successful and failed replications.

\begin{thebibliography}{6}
\small
\bibitem[Li and Rodr\'{i}guez(2025)]{ncc}
R. Li and A. Rodr\'{i}guez. Neural Conformal Control for Time Series Forecasting.
\emph{AAAI}, 39(17):18439--18447, 2025.
\bibitem[Tumu et~al.(2026)]{adapt}
R. Tumu, A. Singh, and R. Mangharam. AdaptNC: Adaptive Nonconformity Scores
for Conformal Prediction under Distribution Shift. arXiv:2602.01629v2, 2026.
\bibitem[Jun and Ohn(2026)]{retro}
J. Jun and I. Ohn. Online conformal inference with retrospective adjustment
for faster adaptation to distribution shift. \emph{Pattern Recognition},
180:114406, 2026.
\bibitem[Bhatnagar et~al.(2023)]{saocp}
A. Bhatnagar, H. Wang, C. Xiong, and Y. Bai.
Improved Online Conformal Prediction via Strongly Adaptive Online Learning.
\emph{ICML}, PMLR 202:2337--2363, 2023.
\bibitem[D'Amour et~al.(2022)]{damour}
A. D'Amour et al. Underspecification Presents Challenges for Credibility in
Modern Machine Learning. \emph{JMLR}, 23(226):1--61, 2022.
\bibitem[Reisenhofer et~al.(2022)]{harnet}
R. Reisenhofer, X. Bayer, and N. Hautsch.
HARNet: A Convolutional Neural Network for Realized Volatility Forecasting.
arXiv:2205.07719, 2022.
\end{thebibliography}
\end{document}
